\documentclass{article}
\usepackage{amsmath}
\usepackage{amssymb}
\usepackage{mathtools}

\begin{document}
\section{Copy paste útiles}

\begin{enumerate}
\item{$\mathbb{Z}$}
\item $x \in l$ 
\item $\wedge$
\item $\vee$
\item $\rightarrow$
\item $x \equiv 0 \mod{y}$
\item $\sum_{i=1}^{n} a_i$

\end{enumerate}

\section{Guia 4}

\textbf{Respuestas}\newline

\textbf{Ejercicio 1.}\newline

\textbf{Ejercicio 2.}\newline

\textbf{Ejercicio 3.}: Especificar e implementar la función esDivisible :: Integer ->Integer -> Bool que dados dos números naturales determinar si el primero es divisible por el segundo. No esta permitido utilizar las funciones mod ni div \newline

problema esDivisible($x: \mathbf{Z}, y: \mathbf{Z}$): Bool \{

    requiere: \{$x>0 \wedge y>0$\}
    
    asegura: \{$x \equiv 0 \mod{y}$\}
    
\} \newline

\textbf{Ejercicio 4.} Especificar e implementar la función sumaImpares::Integer->Integer que dado $n \in \mathbf{N}$ sume los primeros n números impares. Por ejemplo: sumaImpares 3 -> 1+3+5 = 9 \newline

problema sumaImpares($x: \mathbf{Z}$): $\mathbf{Z}$ \{
    
    requiere: \{$x>0$\}

    asegura: \{$res = \sum_{i=1}^{x} (2*n-1)$\}
\}

\textbf{Ejercicio 7}
\end{document}



